\RequirePackage{etex}

% Vi benytter 2 sidet, a4 størrelse og dansk sprog. 
\documentclass[11pt,a4paper,oneside,danish]{memoir}
\usepackage[paper=A4,pagesize]{typearea}
 \usepackage[utf8]{inputenc} % danske tegn

% Choose language by uncommenting danish or english
\usepackage[danish]{babel} % dansk sprogpakke
% \usepackage[english]{babel}


%% Pakker til fonts, farver, referencer og ligenende er placeret relevant sted i dokumentet.
     
     %% FONT
     \usepackage[default,osfigures,scale=0.95]{opensans} % Alternatively use the option 'defaultsans' instead of 'default' to replace the sans serif font only.
     \usepackage[T1]{fontenc}
     
     % Common packages 
     \usepackage{footnote} % Package used to controle footnote(specially in tables etc.)
     \usepackage{amsmath}
     \usepackage{slantsc}
     \usepackage{multirow} % Merge table rows
     \usepackage{longtable} % Table that can span multiple pages
     \usepackage{geometry}
     \usepackage[section]{placeins} % \FloatBarrier
     \usepackage[compatibility=false]{caption}
     \captionsetup{font=small,labelfont=bf,textfont=it}
     \usepackage{graphicx}
     \usepackage{afterpage}
     \usepackage[draft]{fixme} %hvor 'options' bl.a. inkluderer final, draft, danish, inline og margin
     \usepackage{rotating}
     \usepackage{eso-pic} % til inds�ttelse af forside.
     \usepackage{rotating} % rotate stuff(EVERYTHING!11)

	% % MATH STUFF 
    %\usepackage{algorithmic}
    %\usepackage{harpoon} %Tilf�jer harpun-pile over st�rrelser i math, m.m.
    %\usepackage{pst-plot} % For axes   
    %\usepackage{breqn}
    %\usepackage{fourier}  %................... Roman+math - Utopia
    
  
    %\usepackage{nomentbl} % http://ctan.mackichan.com/macros/latex/contrib/nomentbl/nomentbl.pdf

    %\usepackage{psfrag} % insæt tag, mest brugt i forhold til grafer

%%% DIVERSE CUSTOMS PAKKET NED I FILER FOR OVERBLIK
   % \input{styles/HaxTilPreamble.tex}
   % \input{styles/CustomFunctions.tex}
%	\input{Styles/code.tex} 
    \usepackage{todonotes}


% new commands for individually colored todo 
\newcommand{\todosander}[2]{
\todo[#1, color=green!40]{#2 \\- Sander}
}
\newcommand{\todosimon}[2]{
\todo[#1, color=red!40]{#2 \\- Simon}
}
\newcommand{\todorasmus}[2]{
\todo[#1, color=yellow!40]{#2 \\- Rasmus}
}
\newcommand{\todokristoffer}[2]{
\todo[#1, color=purple!40]{#2 \\- Simon}
}


% Color package. Use with listings package and use for color a table
%\usepackage[usenames,dvipsnames,svgnames,table]{xcolor}

    %%%%%%%%%%%%%%%%%%%%%%%%%%%%%%%%%%%%%%%%%%%%%%%%%%%%%%%%%%%%%%%%%%%%%%%%%%%%%
%                            Custom for funktioner                          %
%                     			      	        		                    %
%%%%%%%%%%%%%%%%%%%%%%%%%%%%%%%%%%%%%%%%%%%%%%%%%%%%%%%%%%%%%%%%%%%%%%%%%%%%%


% Package for wrapping figures in text
\usepackage{wrapfig}

% Set the path where LaTeX looks for pictures.
\graphicspath{{figure/}}

% You need a newsubfloat element to use subcaption
\newsubfloat{figure}

% Command to set caption styles
\captionnamefont{\bfseries\small}
\captiontitlefont{\itshape\small}
\subcaptionlabelfont{\bfseries\small}
\subcaptionfont{\itshape\small}

% When we don't use chapter numbering, 
% the figure numbering are shown like 
% fx 0.3 if we don't use this line.
\counterwithout{figure}{chapter}


% Command \figur{filename}{caption}{label}{width} 
% for inserting a new figure.
\newcommand{\figur}[4]{ 
 \begin{figure}[ht] 
 \centering 
 \includegraphics[width=#4\textwidth]{#1} 
 \caption{#2} 
 \label{#3} 
 \end{figure}
 } 
 
% Command \wrapfigur{filename}{caption}{label}{width}{r/l}
% for inserting a new figure wraped in the text. 
\newcommand{\wrapfigur}[5]{ 
 \begin{wrapfigure}{#5}{#4 \textwidth}
   \begin{center}
    \includegraphics[width=#4\textwidth]{#1}
   \end{center}
   \caption{#2}
   \label{#3} 
 \end{wrapfigure}
}
 
% Command \subfigur{filename1}{filename2}{caption}
% {subcaption1}{subcaption2}{label1}{label2}
% {width1}{width1} for inserting a new figure.
\newcommand{\subfigur}[9]{
\begin{figure}[ht]
\centering
\subbottom[{#4}\label{#6}]%
    {\includegraphics[width=#8\textwidth]{#1}}\hfill
\subbottom[{#5}\label{#7}]%
    {\includegraphics[width=#9\textwidth]{#2}}
\caption{#3}
\label{fig:subfigures}
\end{figure}
}
% Command \truefigure{filename}{caption}{label}
% for inserting a picture without scaling.
\newcommand{\truefigure}[3]{ 
 \begin{figure}[ht] 
 \centering 
 \includegraphics{#1} 
 \caption{#2} 
 \label{#3} 
 \end{figure}
 } 

% Keeps floats in the section. 
% Help control place to put pix
\usepackage[section]{placeins}

% Commands for rightaligning wraped text in a tabular
\newcommand{\rr}{\raggedright}
\newcommand{\tn}{\tabularnewline}

\usepackage{multirow}

%Tabular pakke som Wex har fundet, men som vi ikke bruger alligvel
%\usepackage{tabu}

% Include code pieces from files
% Documentation:
% ftp://ftp.tex.ac.uk/tex-archive/macros/latex/contrib/listings/listings.pdf
\usepackage{listings}


%Package with enumirate command with extra options (used within cells of tabular)
\usepackage[inline]{enumitem}
    
    % add a box using the themecolor. Can be used to wrap around a chapintro for introducing an assignment
    %use \begin{assignment} and end\{assignment} to box it
    %http://texdoc.net/texmf-dist/doc/latex/mdframed/mdframed.pdf
     
    \usepackage[xcolor]{mdframed} 
    \newmdenv[linecolor=ThemeColor]{assignment}
%    \usepackage{Placeins}
    
% % % % % % % % % % % % % % ********% % %Generelt design % % % % % % % % % % % % % % % % % % %
 
%\usepackage[table,dvipsnames]{xcolor} % 'table' er i stedet for \usepackage{colortbl} og 'dvipsnames' er et farveskema hvor man angiver farver med navne

\usepackage{xcolor} % farver kan indsætttes i henhold til RGB med definecolor

%Valg af temafarve % Denne farve bruges i chapterstyle %
\definecolor{ThemeColor}{RGB}{209, 67, 26}

% Valg af font %
\usepackage[sc]{mathpazo}
\linespread{1.05}         % Palatino needs more leading (space between lines)
\usepackage[T1]{fontenc}

% Pakker inkludere overordnede designvalg. Pakkerne er valgt for overskuelighed %
	 %%%%%%%%%%%%%%%%%%%%%%%%%%%%%%%%%%%%%%%%%%%%%%%%%%%%%%%%%%%%%%%%%%%%%%%%%%%%%
%                             Chapter and Section                           %
%                     			      Style	       		                    %
%%%%%%%%%%%%%%%%%%%%%%%%%%%%%%%%%%%%%%%%%%%%%%%%%%%%%%%%%%%%%%%%%%%%%%%%%%%%%


% % % % Redefined so that header will show even on chapter start. % % % %
\makeatletter
    \let\stdchapter\chapter
    \renewcommand*\chapter{%
    \@ifstar{\starchapter}{\@dblarg\nostarchapter}}
    \newcommand*\starchapter[1]{%
        \stdchapter*{#1}
        \thispagestyle{fancy}
        \markboth{\MakeUppercase{#1}}{}
    }
    \def\nostarchapter[#1]#2{%
        \stdchapter[{#1}]{#2}
        \thispagestyle{fancy}
    }
\makeatother




\usepackage{titlesec} % Pakke til opsætning af chapter og sektioner


% % defining new font for the chapternumber. Calling it FuckingHuge 
\newcommand{\FuckingHuge}{% Needs a 'stretchable' font
      \usefont{\encodingdefault}{\rmdefault}{b}{n}%
      \fontsize{35}{45}\normalfont\color{ThemeColor}} %60 80


%%% TAL I OVERSKRIFTER, både chapters og sections. 
    \maxsecnumdepth{subsubsection} %subsubsection er dybdeniveau for numering.
   	\setsecnumdepth{subsubsection}


  
 % % % % CHAPTER OPSÆTNING % % % % % % % % % % % % % % % % % % % % %

%Hang thoose numbers! I margin med tallene for både sections og chapters! En lækker chapter style!
\chapterstyle{hangnum}
 
%Modificering af ovenstående style
%Placering af chapter overskriften på siden i forhold til sidens start %
\renewcommand*{\chapterheadstart}{\vspace*{-20pt}}


% FARVE PÅ CHAPTER % %
%TItel farve og font valg
\renewcommand*{\chaptitlefont}{\normalfont\Huge\color{ThemeColor}}
%Numering farve og font valg
\renewcommand{\chapnumfont}{\normalfont\FuckingHuge\color{ThemeColor}}



% % % % SECTION OPSÆTNING % % % % % % % % % % % % % % % % % % % %
\hangsecnum % vi smidder numering ud i marken

%Vi benytter tema farven til numeringen
\renewcommand\thesection{\color{ThemeColor}\thechapter.\arabic{section}}



    
    


	 \input{Styles/header.tex}
	 
% Setup the margins:
\setlrmarginsandblock{1in}{*}{1.5}
\setulmarginsandblock{3cm}{4cm}{*}

% Setup header and footer: %
% \setheadfoot{headheight}{footskip}
% \setheaderspaces{headdrop}{headsep}{ratio}
\setheadfoot{45pt}{1cm}
\setheaderspaces{1cm}{*}{*}
\checkandfixthelayout
% Se dette pdf for dokmentaton:
% www.tug.dk/_media/foredrag/2003-12-03/slides-src.pdf

 \checkandfixthelayout[nearest]
    \checkthelayout
    \fixthelayout
    \baselineskip \topskip

%BibTex setup
\usepackage[semicolon,authoryear,square]{natbib}

\bibliographystyle{plainnat}
%%%%%%%%%%%%%%%%%%%%%%%%%%%%%%%%%%%%%%%%%%%%%%%%%%%%%%%%%%%%%%%%%%%%%%%%%%%%%%%%%%%%%%%%%%%%%%

% % Smart listing function. Documentation: http://texblog.org/tag/makenomenclature/
\usepackage{nomencl}
\makenomenclature 

% % % % % % % % % % % % % % % % Reference opsætning % % % % % % % % % % % % % % % % % % % % % % %

\usepackage{varioref} % for references, http://www.ctex.org/documents/packages/bibref/varioref.pdf
   
% External reference package, zref makes the references into hyperlinks.
\usepackage{url}
\usepackage{hyperref}
%\usepackage{zref-xr}
\usepackage{xr}

% % % % % % % % % % % % % % % % % % % % % % % % % % % % % % % % % % % % % % % % % % % % % % % % % % %

%\usepackage{showlabels}

% footruleskip is allready defined in memoir, so we have to undefine it, in order to include fancyhdr for our footer and header.
\let\footruleskip\undefined
\usepackage{fancyhdr}% For making fancy headers http://ctan.org/pkg/fancyhdr

\AtBeginDocument{\renewcommand\contentsname{Indholdsfortegnelse}}

% % % % % % % % % % % % % Inkluderinger af ønskede kapitler % % % % % % % % % % % % %

\begin{document}
\newgeometry{left=3cm,bottom=0.1cm}
% Write the title
{
\centerline{\Huge\huge\bfseries\color{ThemeColor}Team 15107} 
\vspace{1em}
\centerline{\large\color{ThemeColor}Aarhus School of Engineering - Aarhus Universitet}
}

\centerline{\large\date{10. April 2015}}

%\color{Black}
%\vspace{8em} % Tilføj nogle tomme linjer
\begin{center}
    \begin{tabular}{ l l p{6cm} }
	 & & \\
	& & \\
    Stud. nr.: 201270115 & Johnny Davidsen Poulsen & \\\hline
    & & \\
    Stud. nr.:  02255 & Kim Rose Faber & \\\hline
    & & \\
    Stud. nr.:  & Kristian Rossen & \\\hline
     \end{tabular}
      	\vspace{50pt}
      	
      \textbf{Repository}\\
     http://gitswat.ase.au.dk/git/15/1/team15107/ex13.git
\end{center}
\thispagestyle{empty} % No header and footer on this page
\restoregeometry
\setcounter{tocdepth}{2}
\tableofcontents

 

% Setup header and footer
\pagestyle{fancy}
\fancyhf{} % Clear all header an footer fields 
\renewcommand{\headrulewidth}{0pt} % Delete headruler
\fancyhead[R]{\includegraphics[height=40pt]{header}} % Insert picture in left side of header
%\fancyhead[R]{\rightmark} % insert chapter title in right side of header
\fancyfoot[LE,RO]{\thepage} % insert pagenumber on right side of even page, left side of odd page 
%\setcounter{page}{1}

\thispagestyle{fancy}
\chapter{Beskrivelse}
\label{chp:description}
\todo[inline]{dets indhold, formaal, konsekvenser, etc.}

\section{Indhold}
todo

\section{Formål}
todo

\section{Konsekvenser}
todo
\thispagestyle{fancy}
\chapter{Sammenligning}
\label{chp:comparison}

\section{Problemet uden IOC}

For at forstå IOC er et godt udgangspunkt problematikkerne der kan opstå hvis man ikke benytter sig af dette designprincip. Nedenstående eksempel består af en Customer klasse som indeholder et Address objekt. Problemet er den tætte kobling mellem Customer og Address - hvis Address klassen på noget tidspunkt ændrer sig, vil dette også have indflydelse på kompileringen af Customer klassen.

\figur{Customer1.png}{Customer og Address problemet}{fig:CustomerOgAddress}{0.4}

Det primære problem, som skaber den tætte kobling er, at Customer klassen skal oprette et Address objekt.
Endvidere er Customer klassen koblet direkte op mod Address klassen, så tilføjelse af eks. HomeAddress eller OfficeAddress vil ikke kunne lade sig gøre.

Generelt ses det også, at hvis Customer ikke kan oprette et Address objekt vil det ikke være muligt, at konstruere et Customer objekt.


\section{Generel løsning af problemet}

Problemerne opstår generelt når oprettelsen af et objekt foregår inde i en klassen. Det vil derfor, i Customer Address eksemplet, være ideelt at flytte oprettelsen af Address objektet fra Customer klassen til en anden klasse. I stedet for at Customer har kontrollen af oprettelse af objekter, gives denne til en anden klassen. Dette er den primære ide bag design princippet Inversion of Control (IOC).
Dette princip er også defineret som Hollywood princippet: "Don't call us we will call you" eller fra Address objektets synsvinkel: "Don't create me I will create my self using some one else".


\section{Specifikke måder at implementere IOC}

IOC er et design princip, som i sig selv ikke løser problemerne. Det er her Dependency Injection (DI) kommer ind i billedet, som er en måde at implementere IOC på. 

\figur{IOC.png}{IOC i relation til DI}{fig:IOCandDI}{0.6}

Der er fire måder hvor på DI kan implementeres: 

\begin{itemize}
	\item Constructor
	\item Getter og Setter
	\item Interface
	\item Service locator
\end{itemize}

Alt afhængig af hvordan DI benyttes i programmet, kan de forskellige DI metoder benyttes. Neden for ses et ekempel med Constructor Dependency Injection. I vores tilfælde vil Address blive oprettet af klienten og blot sendt med som constructor parameter til Customer. Nu har Customer en reference til Address og kan frit benytte sig af denne.

\figur{DI.png}{Customer og Address constructor DI}{fig:constructorDI}{0.5}

\todo[inline]{Sammenligne emnet med andre, tilsvarende emner. Disse maa gerne have vaeret behandlet i undervisningen.}

\section{Factory}

Den faktiske måde DI bliver implementeret på er også forskellig. Det er muligt, at oprette alle de afhængige objekter i en enkelte klasse, eksempelvis main(), og sende dem videre ned i de oprettede objekter. Det primære problem er, at når oprettelsen af objekter sker i en enkelt klasse opstår der alligevel høj kobling. Eksempelvis, når en concrete klasse skal udvides med en ekstra constructor parameter, skal alle overstående klasser implementere denne. Dette kan dog løses ved hjælp af Factory design mønstret. 

\figur{Factory.png}{Factory design pattern}{fig:constructorDI}{0.5}

Her sendes der blot et Factory objekt rundt, som kan opfylde oprettelsen af de enkelte objekters behov. Selv om Factory kan løse mange af problemerne, er der dog stadig argumenter for, at dette ikke er et optimalt design. Det største problem med en Factory er, at den kan ikke genbruges i andre applikationer. Alle valg om oprettelse er med henblik på den nuværende applikation, og det er sjældent muligt, at klasser til en anden applikation. 

\todo {Interface dependent and compile time}


\section{Container}

En Container kan tage ansvaret for at oprette, konfigurere og holde styr på objekterne. Det er muligt, at konfigurere objekterne i Containeren, som også hjælper med at oprette og injecte objekterne i de respektive klasser. 

I tilfældet med Customer og Address, vil Containeren være den, som registrer de to objekter separat og injecter Address objekt ind i Customer klassen. Disse Containers er ofte tilbudt som et framework, der gør det let at implementere og har et højt abstraktionsniveau.
\thispagestyle{fancy}
\chapter{Eksempler}
\label{chp:examples}
\todo[inline]{Eksemplificere emnet – for eksempel vise en implementering af emnet, et eksempel-gennemløb af udviklingsprocessen eller lignende.}


\thispagestyle{fancy}
\chapter{Konklusion}
\label{chp:Konklusion}
Der har gennem projektperioden vist sig et behov for at teste. Flere tests har afsløret væsentlige fejl hvor der eksempelvist blev ittereret gennem lister i forkert rækkefølge eller beregninger blev returneret med værdier som havde forkerte værdier på 12. decimal. Anden itteration gav et design som er meget nemmere at udvidde med lignende krav. Fordelen ved at implementere et design som bygger på SOLID principperne har vist sine fordele. Vi har gennem dette projekt udviklet en bedre forståelse for SOLID og relevansen af unittesting. Integrationenstestene er desværre ikke implementeret, men der er tilgengæld lavet en række test ved hjælp af et eksekverbar projekt hvor en række scenarier er testet manuelt. Dette har givet os en god fornemmelse af at projektet fungere uden at vi dog har automatiseret processen.

%I forbindelse med opgaven har gruppen fået en bedre forståelse for vigtigheden af en godt og testbart design med en klar opdeling af klasserne. Dette gjorde det også forholdsvist smertefrit at håndtere de ændrede krav, der kom sent i forløbet, og det har givet gruppen et indblik i, hvad man kan blive udsat for i et virkeligt projekt. 
%
%Brugen af Jenkins var en motiverende faktor for at lave tests løbende, da man kan følge med i grafens udvikling, og det har haft stor betydning for udviklingsforløbet at projektet. 
%
%Ydermere har gruppen fået et væsentligt bedre erfaring til brugen af testframeworks i form af NUnit, DotCover og NSubstitute. Både i forhold til det rent tekniske, men i samme grad at skabe en god struktur og præcision i de udarbejdede tests.


%\bibliographystyle{plainnat}
%\bibliography{bibliografi}% Selects .bib file AND prints bibliography
\thispagestyle{fancy}

%\listoftodos % Make a list of todo's
\end{document}